% --- Archon physics formalization source begin ---
% archon:physics
% archon:covers PhyXMiniProblems/problem_phyx_mini_0873.lean
% archon:source-report reports/phyx_mini/problem_phyx_mini_0873.source.json

\chapter{Physics problem phyx\_mini\_0873}
\label{ch:PhyXMiniProblems_problem_phyx_mini_0873}

\paragraph{Problem source.}
\#\# Question

What is the capacitance of the two metal spheres?

\#\# Answer choices

- A: A: 0.22nF

- B: B: 1.20nF

- C: C: 0.25nF

- D: D: 0.20nF

\#\# Figure caption (auxiliary; use the image as primary evidence)

The image depicts two circles representing charges. 

- On the left side is a circle with a red plus sign inside, labeled "+20 nC," indicating a positive charge of 20 nanocoulombs.
- On the right side is a circle with a teal minus sign inside, labeled "−20 nC," indicating a negative charge of 20 nanocoulombs.
- Between the two circles is a double-headed green arrow pointing left and right, with text underneath stating "ΔV = 100 V," indicating a potential difference of 100 volts between the charges.

\#\# Dataset metadata

Electrostatics; Physical Model Grounding Reasoning, Multi-Formula Reasoning

\paragraph{Recorded answer/context.}
D: D: 0.20nF

\paragraph{Figure/image path.}
/root/proposal\_for\_physic/science-mango/phyx\_mini\_run/phyx\_data/test\_image/873.png

\paragraph{Formalization target.}
create a compiling Lean file with sorry bodies at `PhyXMiniProblems/problem\_phyx\_mini\_0873.lean`. The Lean declarations must preserve the physical quantities, dimensions or dimensional roles, figure labels, governing-law hypotheses, and final relation expressed by this problem.
Use Mathlib/Physlib names found through LeanExplore where available. If a physics API is missing, introduce faithful local abstractions rather than scalar placeholder aliases.

\begin{theorem}[Physics formalization target]
\label{thm:physics:phyx_mini_0873:target}
\lean{PhyXMiniProblems.ProblemPhyXMini0873.problem_phyx_mini_0873}
\uses{def:physics:phyx-mini-0873:phyxminiproblems-problemphyxmini0873-capacitanceinnanofarads, def:physics:phyx-mini-0873:phyxminiproblems-problemphyxmini0873-twometalspherecapacitorsetup, def:physics:phyx-mini-0873:phyxminiproblems-problemphyxmini0873-matchesproblemandfigurereadouts, def:physics:phyx-mini-0873:phyxminiproblems-problemphyxmini0873-hasphysicalparameters, def:physics:phyx-mini-0873:phyxminiproblems-problemphyxmini0873-satisfiestwoconductorcapacitancelaw, def:physics:phyx-mini-0873:phyxminiproblems-problemphyxmini0873-recordeddatasetanswer, def:physics:phyx-mini-0873:phyxminiproblems-problemphyxmini0873-answermatchescapacitance}
The assigned autoformalize agent should translate this physics problem into Lean declarations in the covered file, with theorem and lemma proof bodies written as `by sorry`.
\end{theorem}
\begin{proof}
This is an autoformalization task, not a proof task. Produce statements that can later be proved by the physics prover without weakening the physical model.
\end{proof}

% --- Archon declaration topology begin ---
\section{Declaration topology}
\begin{definition}[potential Difference Dimension]
  \label{def:physics:phyx-mini-0873:phyxminiproblems-problemphyxmini0873-potentialdifferencedimension}
  \lean{PhyXMiniProblems.ProblemPhyXMini0873.potentialDifferenceDimension}
  The dimension M L² T⁻² C⁻¹ of electric potential difference
\end{definition}
\begin{proof}
  This is the declared model definition of potential Difference Dimension.
\end{proof}
\begin{definition}[capacitance Dimension]
  \label{def:physics:phyx-mini-0873:phyxminiproblems-problemphyxmini0873-capacitancedimension}
  \lean{PhyXMiniProblems.ProblemPhyXMini0873.capacitanceDimension}
  \uses{def:physics:phyx-mini-0873:phyxminiproblems-problemphyxmini0873-potentialdifferencedimension}
  Capacitance has dimension charge divided by potential difference
\end{definition}
\begin{proof}
  This is the declared model definition of capacitance Dimension.
\end{proof}
\begin{definition}[Signed Charge Quantity]
  \label{def:physics:phyx-mini-0873:phyxminiproblems-problemphyxmini0873-signedchargequantity}
  \lean{PhyXMiniProblems.ProblemPhyXMini0873.SignedChargeQuantity}
  A signed, unit-independent physical electric charge
\end{definition}
\begin{proof}
  This is the declared model definition of Signed Charge Quantity.
\end{proof}
\begin{definition}[Potential Difference Quantity]
  \label{def:physics:phyx-mini-0873:phyxminiproblems-problemphyxmini0873-potentialdifferencequantity}
  \lean{PhyXMiniProblems.ProblemPhyXMini0873.PotentialDifferenceQuantity}
  \uses{def:physics:phyx-mini-0873:phyxminiproblems-problemphyxmini0873-potentialdifferencedimension}
  A nonnegative, unit-independent potential-difference magnitude
\end{definition}
\begin{proof}
  This is the declared model definition of Potential Difference Quantity.
\end{proof}
\begin{definition}[Capacitance Quantity]
  \label{def:physics:phyx-mini-0873:phyxminiproblems-problemphyxmini0873-capacitancequantity}
  \lean{PhyXMiniProblems.ProblemPhyXMini0873.CapacitanceQuantity}
  \uses{def:physics:phyx-mini-0873:phyxminiproblems-problemphyxmini0873-capacitancedimension}
  A nonnegative, unit-independent physical capacitance
\end{definition}
\begin{proof}
  This is the declared model definition of Capacitance Quantity.
\end{proof}
\begin{definition}[charge In Coulombs]
  \label{def:physics:phyx-mini-0873:phyxminiproblems-problemphyxmini0873-chargeincoulombs}
  \lean{PhyXMiniProblems.ProblemPhyXMini0873.chargeInCoulombs}
  \uses{def:physics:phyx-mini-0873:phyxminiproblems-problemphyxmini0873-signedchargequantity}
  Coherent-SI readout of a signed electric charge, in coulombs
\end{definition}
\begin{proof}
  This is the declared model definition of charge In Coulombs.
\end{proof}
\begin{definition}[charge In Nanocoulombs]
  \label{def:physics:phyx-mini-0873:phyxminiproblems-problemphyxmini0873-chargeinnanocoulombs}
  \lean{PhyXMiniProblems.ProblemPhyXMini0873.chargeInNanocoulombs}
  \uses{def:physics:phyx-mini-0873:phyxminiproblems-problemphyxmini0873-signedchargequantity, def:physics:phyx-mini-0873:phyxminiproblems-problemphyxmini0873-chargeincoulombs}
  Nanocoulomb readout used by the two charge labels in the figure
\end{definition}
\begin{proof}
  This is the declared model definition of charge In Nanocoulombs.
\end{proof}
\begin{definition}[potential Difference In Volts]
  \label{def:physics:phyx-mini-0873:phyxminiproblems-problemphyxmini0873-potentialdifferenceinvolts}
  \lean{PhyXMiniProblems.ProblemPhyXMini0873.potentialDifferenceInVolts}
  \uses{def:physics:phyx-mini-0873:phyxminiproblems-problemphyxmini0873-potentialdifferencequantity}
  Coherent-SI readout of a potential difference, in volts
\end{definition}
\begin{proof}
  This is the declared model definition of potential Difference In Volts.
\end{proof}
\begin{definition}[capacitance In Farads]
  \label{def:physics:phyx-mini-0873:phyxminiproblems-problemphyxmini0873-capacitanceinfarads}
  \lean{PhyXMiniProblems.ProblemPhyXMini0873.capacitanceInFarads}
  \uses{def:physics:phyx-mini-0873:phyxminiproblems-problemphyxmini0873-capacitancequantity}
  Coherent-SI readout of a capacitance, in farads
\end{definition}
\begin{proof}
  This is the declared model definition of capacitance In Farads.
\end{proof}
\begin{definition}[capacitance In Nanofarads]
  \label{def:physics:phyx-mini-0873:phyxminiproblems-problemphyxmini0873-capacitanceinnanofarads}
  \lean{PhyXMiniProblems.ProblemPhyXMini0873.capacitanceInNanofarads}
  \uses{def:physics:phyx-mini-0873:phyxminiproblems-problemphyxmini0873-capacitancequantity, def:physics:phyx-mini-0873:phyxminiproblems-problemphyxmini0873-capacitanceinfarads}
  Nanofarad readout used by the four displayed answer choices
\end{definition}
\begin{proof}
  This is the declared model definition of capacitance In Nanofarads.
\end{proof}
\begin{definition}[Sphere Label]
  \label{def:physics:phyx-mini-0873:phyxminiproblems-problemphyxmini0873-spherelabel}
  \lean{PhyXMiniProblems.ProblemPhyXMini0873.SphereLabel}
  The two physical spheres, named by their positions in the figure
\end{definition}
\begin{proof}
  This is the declared model definition of Sphere Label.
\end{proof}
\begin{definition}[Horizontal Position]
  \label{def:physics:phyx-mini-0873:phyxminiproblems-problemphyxmini0873-horizontalposition}
  \lean{PhyXMiniProblems.ProblemPhyXMini0873.HorizontalPosition}
  Horizontal positions available to the two sphere drawings
\end{definition}
\begin{proof}
  This is the declared model definition of Horizontal Position.
\end{proof}
\begin{definition}[Body Material]
  \label{def:physics:phyx-mini-0873:phyxminiproblems-problemphyxmini0873-bodymaterial}
  \lean{PhyXMiniProblems.ProblemPhyXMini0873.BodyMaterial}
  Material classification needed to state that both bodies are metal
\end{definition}
\begin{proof}
  This is the declared model definition of Body Material.
\end{proof}
\begin{definition}[Body Shape]
  \label{def:physics:phyx-mini-0873:phyxminiproblems-problemphyxmini0873-bodyshape}
  \lean{PhyXMiniProblems.ProblemPhyXMini0873.BodyShape}
  Shape classification needed to preserve the spherical geometry
\end{definition}
\begin{proof}
  This is the declared model definition of Body Shape.
\end{proof}
\begin{definition}[Figure Charge Sign]
  \label{def:physics:phyx-mini-0873:phyxminiproblems-problemphyxmini0873-figurechargesign}
  \lean{PhyXMiniProblems.ProblemPhyXMini0873.FigureChargeSign}
  Sign glyph drawn at the centre of a sphere in the supplied image
\end{definition}
\begin{proof}
  This is the declared model definition of Figure Charge Sign.
\end{proof}
\begin{definition}[expected Horizontal Position]
  \label{def:physics:phyx-mini-0873:phyxminiproblems-problemphyxmini0873-expectedhorizontalposition}
  \lean{PhyXMiniProblems.ProblemPhyXMini0873.expectedHorizontalPosition}
  \uses{def:physics:phyx-mini-0873:phyxminiproblems-problemphyxmini0873-spherelabel, def:physics:phyx-mini-0873:phyxminiproblems-problemphyxmini0873-horizontalposition}
  Horizontal position expected for each named sphere
\end{definition}
\begin{proof}
  This is the declared model definition of expected Horizontal Position.
\end{proof}
\begin{definition}[expected Charge Sign]
  \label{def:physics:phyx-mini-0873:phyxminiproblems-problemphyxmini0873-expectedchargesign}
  \lean{PhyXMiniProblems.ProblemPhyXMini0873.expectedChargeSign}
  \uses{def:physics:phyx-mini-0873:phyxminiproblems-problemphyxmini0873-spherelabel, def:physics:phyx-mini-0873:phyxminiproblems-problemphyxmini0873-figurechargesign}
  Sign glyph expected inside each sphere
\end{definition}
\begin{proof}
  This is the declared model definition of expected Charge Sign.
\end{proof}
\begin{definition}[expected Printed Charge In Nanocoulombs]
  \label{def:physics:phyx-mini-0873:phyxminiproblems-problemphyxmini0873-expectedprintedchargeinnanocoulombs}
  \lean{PhyXMiniProblems.ProblemPhyXMini0873.expectedPrintedChargeInNanocoulombs}
  \uses{def:physics:phyx-mini-0873:phyxminiproblems-problemphyxmini0873-spherelabel}
  Signed nanocoulomb value printed above each sphere
\end{definition}
\begin{proof}
  This is the declared model definition of expected Printed Charge In Nanocoulombs.
\end{proof}
\begin{definition}[Two Sphere Figure]
  \label{def:physics:phyx-mini-0873:phyxminiproblems-problemphyxmini0873-twospherefigure}
  \lean{PhyXMiniProblems.ProblemPhyXMini0873.TwoSphereFigure}
  \uses{def:physics:phyx-mini-0873:phyxminiproblems-problemphyxmini0873-potentialdifferencequantity, def:physics:phyx-mini-0873:phyxminiproblems-problemphyxmini0873-spherelabel, def:physics:phyx-mini-0873:phyxminiproblems-problemphyxmini0873-horizontalposition, def:physics:phyx-mini-0873:phyxminiproblems-problemphyxmini0873-figurechargesign}
  Literal presentation data transcribed from image 873.png
\end{definition}
\begin{proof}
  This is the declared model definition of Two Sphere Figure.
\end{proof}
\begin{definition}[Two Metal Sphere Capacitor Setup]
  \label{def:physics:phyx-mini-0873:phyxminiproblems-problemphyxmini0873-twometalspherecapacitorsetup}
  \lean{PhyXMiniProblems.ProblemPhyXMini0873.TwoMetalSphereCapacitorSetup}
  \uses{def:physics:phyx-mini-0873:phyxminiproblems-problemphyxmini0873-signedchargequantity, def:physics:phyx-mini-0873:phyxminiproblems-problemphyxmini0873-potentialdifferencequantity, def:physics:phyx-mini-0873:phyxminiproblems-problemphyxmini0873-capacitancequantity, def:physics:phyx-mini-0873:phyxminiproblems-problemphyxmini0873-spherelabel, def:physics:phyx-mini-0873:phyxminiproblems-problemphyxmini0873-bodymaterial, def:physics:phyx-mini-0873:phyxminiproblems-problemphyxmini0873-bodyshape, def:physics:phyx-mini-0873:phyxminiproblems-problemphyxmini0873-twospherefigure}
  The independent physical bodies and electrical observables. In particular, capacitance is not defined from the charge and voltage labels or from an answer choice; it is constrained only by the governing law below
\end{definition}
\begin{proof}
  This is the declared model definition of Two Metal Sphere Capacitor Setup.
\end{proof}
\begin{definition}[Matches Problem And Figure Readouts]
  \label{def:physics:phyx-mini-0873:phyxminiproblems-problemphyxmini0873-matchesproblemandfigurereadouts}
  \lean{PhyXMiniProblems.ProblemPhyXMini0873.MatchesProblemAndFigureReadouts}
  \uses{def:physics:phyx-mini-0873:phyxminiproblems-problemphyxmini0873-chargeinnanocoulombs, def:physics:phyx-mini-0873:phyxminiproblems-problemphyxmini0873-potentialdifferenceinvolts, def:physics:phyx-mini-0873:phyxminiproblems-problemphyxmini0873-expectedhorizontalposition, def:physics:phyx-mini-0873:phyxminiproblems-problemphyxmini0873-expectedchargesign, def:physics:phyx-mini-0873:phyxminiproblems-problemphyxmini0873-expectedprintedchargeinnanocoulombs, def:physics:phyx-mini-0873:phyxminiproblems-problemphyxmini0873-twometalspherecapacitorsetup}
  The material, geometry, and numerical labels supplied by the question and its primary image. No capacitance value or answer label occurs in these data
\end{definition}
\begin{proof}
  This is the declared model definition of Matches Problem And Figure Readouts.
\end{proof}
\begin{definition}[Has Physical Parameters]
  \label{def:physics:phyx-mini-0873:phyxminiproblems-problemphyxmini0873-hasphysicalparameters}
  \lean{PhyXMiniProblems.ProblemPhyXMini0873.HasPhysicalParameters}
  \uses{def:physics:phyx-mini-0873:phyxminiproblems-problemphyxmini0873-chargeincoulombs, def:physics:phyx-mini-0873:phyxminiproblems-problemphyxmini0873-potentialdifferenceinvolts, def:physics:phyx-mini-0873:phyxminiproblems-problemphyxmini0873-capacitanceinfarads, def:physics:phyx-mini-0873:phyxminiproblems-problemphyxmini0873-twometalspherecapacitorsetup}
  Positivity and sign conditions selecting the physical configuration
\end{definition}
\begin{proof}
  This is the declared model definition of Has Physical Parameters.
\end{proof}
\begin{definition}[Satisfies Two Conductor Capacitance Law]
  \label{def:physics:phyx-mini-0873:phyxminiproblems-problemphyxmini0873-satisfiestwoconductorcapacitancelaw}
  \lean{PhyXMiniProblems.ProblemPhyXMini0873.SatisfiesTwoConductorCapacitanceLaw}
  \uses{def:physics:phyx-mini-0873:phyxminiproblems-problemphyxmini0873-chargeincoulombs, def:physics:phyx-mini-0873:phyxminiproblems-problemphyxmini0873-potentialdifferenceinvolts, def:physics:phyx-mini-0873:phyxminiproblems-problemphyxmini0873-capacitanceinfarads, def:physics:phyx-mini-0873:phyxminiproblems-problemphyxmini0873-twometalspherecapacitorsetup}
  For a two-conductor capacitor, the conductors carry equal and opposite charge and the magnitude on either conductor satisfies |Q| = C ΔV. This is the general capacitor law in coherent SI readouts and contains no numerical capacitance from the current question
\end{definition}
\begin{proof}
  This is the declared model definition of Satisfies Two Conductor Capacitance Law.
\end{proof}
\begin{lemma}[capacitance eq charge Magnitude div potential Difference]
  \label{lem:physics:phyx-mini-0873:phyxminiproblems-problemphyxmini0873-capacitance-eq-chargemagnitude-div-potentialdifference}
  \lean{PhyXMiniProblems.ProblemPhyXMini0873.capacitance_eq_chargeMagnitude_div_potentialDifference}
  \uses{def:physics:phyx-mini-0873:phyxminiproblems-problemphyxmini0873-chargeincoulombs, def:physics:phyx-mini-0873:phyxminiproblems-problemphyxmini0873-potentialdifferenceinvolts, def:physics:phyx-mini-0873:phyxminiproblems-problemphyxmini0873-capacitanceinfarads, def:physics:phyx-mini-0873:phyxminiproblems-problemphyxmini0873-twometalspherecapacitorsetup, def:physics:phyx-mini-0873:phyxminiproblems-problemphyxmini0873-hasphysicalparameters, def:physics:phyx-mini-0873:phyxminiproblems-problemphyxmini0873-satisfiestwoconductorcapacitancelaw}
  The capacitor law can be rearranged to C = |Q| / ΔV
\end{lemma}
\begin{proof}
  The conclusion follows from the governing relations and figure readouts named in the statement of capacitance eq charge Magnitude div potential Difference.
\end{proof}
\begin{lemma}[capacitance eq one Fifth nanofarad]
  \label{lem:physics:phyx-mini-0873:phyxminiproblems-problemphyxmini0873-capacitance-eq-onefifth-nanofarad}
  \lean{PhyXMiniProblems.ProblemPhyXMini0873.capacitance_eq_oneFifth_nanofarad}
  \uses{def:physics:phyx-mini-0873:phyxminiproblems-problemphyxmini0873-capacitanceinnanofarads, def:physics:phyx-mini-0873:phyxminiproblems-problemphyxmini0873-twometalspherecapacitorsetup, def:physics:phyx-mini-0873:phyxminiproblems-problemphyxmini0873-matchesproblemandfigurereadouts, def:physics:phyx-mini-0873:phyxminiproblems-problemphyxmini0873-hasphysicalparameters, def:physics:phyx-mini-0873:phyxminiproblems-problemphyxmini0873-satisfiestwoconductorcapacitancelaw}
  Substituting |Q| = 20 nC and ΔV = 100 V gives C = 2 * 10⁻¹⁰ F = 0.20 nF
\end{lemma}
\begin{proof}
  The conclusion follows from the governing relations and figure readouts named in the statement of capacitance eq one Fifth nanofarad.
\end{proof}
\begin{definition}[Answer Choice]
  \label{def:physics:phyx-mini-0873:phyxminiproblems-problemphyxmini0873-answerchoice}
  \lean{PhyXMiniProblems.ProblemPhyXMini0873.AnswerChoice}
  Labels of the four capacitance choices printed in the question
\end{definition}
\begin{proof}
  This is the declared model definition of Answer Choice.
\end{proof}
\begin{definition}[answer Capacitance In Nanofarads]
  \label{def:physics:phyx-mini-0873:phyxminiproblems-problemphyxmini0873-answercapacitanceinnanofarads}
  \lean{PhyXMiniProblems.ProblemPhyXMini0873.answerCapacitanceInNanofarads}
  \uses{def:physics:phyx-mini-0873:phyxminiproblems-problemphyxmini0873-answerchoice}
  Capacitance in nanofarads printed beside each answer choice
\end{definition}
\begin{proof}
  This is the declared model definition of answer Capacitance In Nanofarads.
\end{proof}
\begin{definition}[recorded Dataset Answer]
  \label{def:physics:phyx-mini-0873:phyxminiproblems-problemphyxmini0873-recordeddatasetanswer}
  \lean{PhyXMiniProblems.ProblemPhyXMini0873.recordedDatasetAnswer}
  \uses{def:physics:phyx-mini-0873:phyxminiproblems-problemphyxmini0873-answerchoice}
  The answer label recorded in the source dataset
\end{definition}
\begin{proof}
  This is the declared model definition of recorded Dataset Answer.
\end{proof}
\begin{definition}[Answer Matches Capacitance]
  \label{def:physics:phyx-mini-0873:phyxminiproblems-problemphyxmini0873-answermatchescapacitance}
  \lean{PhyXMiniProblems.ProblemPhyXMini0873.AnswerMatchesCapacitance}
  \uses{def:physics:phyx-mini-0873:phyxminiproblems-problemphyxmini0873-capacitanceinnanofarads, def:physics:phyx-mini-0873:phyxminiproblems-problemphyxmini0873-twometalspherecapacitorsetup, def:physics:phyx-mini-0873:phyxminiproblems-problemphyxmini0873-answerchoice, def:physics:phyx-mini-0873:phyxminiproblems-problemphyxmini0873-answercapacitanceinnanofarads}
  A displayed choice matches the independently modeled capacitance
\end{definition}
\begin{proof}
  This is the declared model definition of Answer Matches Capacitance.
\end{proof}
% --- Archon declaration topology end ---
% --- Archon physics formalization source end ---
